\documentclass[11pt]{article}
\usepackage[margin=1in]{geometry}
\usepackage{booktabs,hyperref,amsmath}
\newcommand{\todo}[1]{\textbf{[TODO: #1]}}
\newcommand{\unread}{$^{\dagger}$}
\title{How Much of an RDF Engine's Ranking Is Its Index Layout?\\\large An Experiments \& Analysis Study (draft)}
\author{TripleLight authors}
\date{Draft, 2026-11-22. Sections 4 and 7 await the dedicated-machine measurements.}
\begin{document}
\maketitle

\noindent\textbf{Status of this draft.}
Every number cites its evidence file (\texttt{docs/...} in the repository). Papers marked with~\unread\ were not read
in full by the author; they are cited from their abstract or from a source we read that describes them, and must be
re-checked before submission. Section~4 holds only the registered predictions and empty tables.

\begin{abstract}
RDF engines are routinely ranked on LUBM, WatDiv and BSBM. Engines differ in which triple-index orders they store, and
a ranking therefore mixes the quality of the executor with the choice of indexes. We ask how much of today's ranking is
the index layout. We rank engines twice on each suite: once in their default configuration and once with the index orders
aligned across systems, and we register in advance that the claim fails if the two rankings agree at Kendall's
$\tau \ge 0.8$ on every suite. As a storage arm we study a store that keeps four of the six orders and derives the other
two exactly. \todo{one sentence with the result once Section 4 is measured}
\end{abstract}

\section{Introduction}
When two RDF engines are compared on a benchmark, the table reports one number per system and suite. It does not usually
report which of the six triple orders (SPO, SOP, PSO, POS, OSP, OPS) each system keeps, although the choice decides
which access paths are available: a pattern whose bound positions are not a prefix of any stored order falls back to a
scan or a filter. Hexastore\unread\ stores all six orders; other engines store fewer. The question of this paper is whether
that choice moves the \emph{ranking}, not only the absolute times.

Qualitative evidence that physical layout dominates RDF performance is old: Sidirourgos et al.\ found in 2008 that
clustering a triple store by PSO instead of SPO changed performance by factors of two to five, and that a triple store with
the right clustered indexes could beat vertical partitioning on a row store~\cite{sidirourgos2008}. What is missing is the
corresponding question for current engines: if systems are re-ranked with their index layouts aligned, does the order
change?

\paragraph{Contributions.}
\begin{enumerate}
  \item \textbf{A ranking test with a registered kill rule.} For each of LUBM, WatDiv and BSBM we compare the ranking of
  systems in their default configuration with the ranking under index-aligned configurations, using Kendall's $\tau$. The
  claim that layout matters fails if $\tau \ge 0.8$ on every suite. The rule, the order set of every arm and the predicted
  direction are fixed before any result exists (Section~4.0).
  \item \textbf{A layout-aligned protocol} that records each system's index inventory, pins versions, applies one answer gate
  to every reported time, and reports wrong, incomplete and timed-out runs as such instead of timing them
  (Section~3).
  \item \textbf{A storage arm.} A store that keeps SPO, PSO, POS and OSP and derives SOP and OPS exactly is
  36\% smaller on LUBM-1000 and 17\% smaller on BSBM-100M than the same store with six orders
  (\texttt{docs/build-compare-orders-v1}). Whether query time changes is re-measured on the dedicated machine (Section~5).
\end{enumerate}
The artifact (raw receipts, scripts, reproduction commands) is described in the appendix. If the ranking does not change under
alignment, the paper reduces to the storage arm and an engineering report on loading; we say so in advance in Section~8.

\section{Related Work}
\paragraph{Index orders in RDF stores.} RDF-3X\unread\ and Hexastore\unread\ keep all or most of the six permutations;
QLever\unread\ offers a two-order mode; The Ring\unread\ answers ordered access over all six orders from one structure.
Our storage arm sits between: four ordinary B-tree-like orders and exact derivation of the remaining two. Where the prior systems are
cited, we have read abstracts or secondary descriptions only (marked~\unread).

\paragraph{Layout versus store.} Sidirourgos et al.\ reexamined column-store RDF storage and found that the advantage
reported for vertical partitioning depends on how the triple store is clustered~\cite{sidirourgos2008}. That study
re-evaluated the vertical-partitioning proposal of Abadi et al.\ on MonetDB/SQL and PostgreSQL; it compared storage schemes, not engines against one another.
We ask the ranking question for current engines.

\paragraph{Benchmark data and structure.} Duan et al.\ showed that the structuredness of LUBM, BSBM and SP$^2$Bench data is
nearly fixed and differs from that of real RDF datasets such as DBpedia, so conclusions drawn from benchmark runs may not transfer
to real data~\cite{duan2011}. This is a threat to our own results (Section~8): our suites are the benchmark datasets they analyse,
and we report them as benchmark results only.

\paragraph{Methodology.} \todo{SPARQL log studies (Bonifati et al.\unread, Arias et al.\unread) and the real-query-log coverage of the derived orders; read before writing.}

\section{Methodology}
\paragraph{Systems.} Versions are pinned in \texttt{docs/eval-plan/PLAN.md} (frozen plan). Server cohort (HTTP): TripleLight,
QLever, Virtuoso~7, Jena TDB2/Fuseki, Oxigraph, RDF4J (native). Embedded cohort (language API): TripleLight bindings,
pyoxigraph, Blazegraph, and DuckDB as a SQL baseline. The embedded and the HTTP cohort are compared within themselves and
never in one figure.

\paragraph{Index inventory.} For every system we record the orders it stores in its measured configuration
(\texttt{docs/embedded-store-inventory.md} and its receipt supplement). Examples: RDF4J's measured profile stores
\texttt{spoc, posc, ospc} (three orders); DuckDB is run as a single table and as six sorted copies, fixed and routed, as
contrasts \emph{within} DuckDB only; TripleLight is run with four and with six orders; QLever with its default and its two-order mode.

\paragraph{Datasets and queries.} LUBM-1000, WatDiv-100M and BSBM-100M, with the benchmark query suites and the analytical
(BI and aggregate) families of the plan. \todo{final scales after the machine is rented.}

\paragraph{Protocol.} One system at a time on one dedicated Linux machine. Cold runs drop the page cache before a fresh
process; warm runs do one warm-up pass and then five measured runs; median and range are reported, geometric mean per suite.
Timeout is 300\,s per query. Every answer is checked against a frozen Jena reference. Each result falls in one class: pass,
value-compatible, mismatch, incomplete, timeout or error; only passing and value-compatible runs are timed, and each figure
names the gate it uses. Variants are compared with ABBA ordering. Equal CPU and memory budgets are enforced by cgroup.

\paragraph{Ranking test.} For every suite we rank the systems of a cohort by geometric-mean time (default configuration), re-rank
under the aligned configurations, and compute Kendall's $\tau$ between the two rankings. Cells where a system cannot answer a
suite are marked unsupported and left out of the aligned ranking, which then covers fewer systems or suites
(Section~8).

\section{Results: Rankings}
\subsection{Registered predictions (written before any measurement)}
\begin{itemize}
  \item \textbf{Quantity:} Kendall's $\tau$ between the default-configuration ranking and the aligned ranking, per suite and cohort.
  \item \textbf{Arms and order sets:}
  TripleLight: \{SPO, PSO, POS, OSP\} and the six-order store; QLever: default and two-order mode (its two-order mode may not
  answer variable-predicate patterns); RDF4J: \{spoc, posc, ospc\}; DuckDB: single table, six sorted copies fixed and routed;
  other systems: as recorded in the inventory. \todo{freeze the exact list from the inventory before the first run.}
  \item \textbf{Kill rule:} the claim that index layout changes the ranking fails if $\tau \ge 0.8$ on every suite.
  \item \textbf{Predicted direction:} \todo{owner of the prediction (Kaiming/Jianlin) to supply the registered direction; the writer does not set it.}
\end{itemize}
\subsection{Results}
\begin{tabular}{lccc}\toprule Suite & Cohort & $\tau$ & Ranks that change\\\midrule
LUBM & embedded & \todo{} & \todo{}\\ LUBM & HTTP & \todo{} & \todo{}\\
WatDiv & embedded & \todo{} & \todo{}\\ WatDiv & HTTP & \todo{} & \todo{}\\
BSBM & embedded & \todo{} & \todo{}\\ BSBM & HTTP & \todo{} & \todo{}\\\bottomrule\end{tabular}

\section{Storage Arm: Four Orders Instead of Six}
TripleLight's default store keeps four orders (SPO, PSO, POS, OSP); SOP and OPS are derived exactly from them at query time.
Dropping the two orders makes the store smaller: on LUBM-1000, from 4.97 to 3.16\,GB ($-36\%$), and on BSBM-100M from 7.94 to
6.56\,GB ($-17\%$) (\texttt{docs/build-compare-orders-v1}). Probe queries over the benchmark suites returned identical
answers in both layouts. The cost of the derived path in query time is \emph{re-measured on the dedicated machine}; no
laptop timing is claimed here. \todo{table: query time, four vs six orders, per suite, with the answer gate.}
How often real queries use a derived order is open and gates whether materializing derived orders on demand is worth it
(\todo{coverage trace on real query logs, task~\#98}). A store written in the four-order format cannot be opened by binaries
that predate its format, and there is no downgrade path (\texttt{docs/release-notes/v0.1.0.md}).

\section{Build and Load}
This section supports the study; it is not a contribution. A four-order build of LUBM-1000 took 654.8\,s. Shard parsing and
sorting took 38.7\% of it, the dictionary-side stages 35\%, and the order builds 26.4\% (\texttt{docs/build-profile-v1}).
In a microbenchmark, fixed-width external sort rows were 4--5$\times$ faster than the current variable-length rows
(\texttt{docs/build-profile-v1}); this is the basis of the fixed-width sorter in the current build. The terms in the data repeat
heavily (7.9\% distinct in LUBM-1000 and 9.3\% per shard in BSBM-100M; \texttt{docs/parse-phase-v1}), which makes
deduplication before sorting attractive. QLever's index builder already deduplicates before sorting, so we treat this as
engineering that follows an existing design, not as a new method.

\section{Ablations}
\todo{Count shortcut, cache budget (\texttt{docs/cache-budget-sweep-v1}), result windowing; reported as ablations.}

\section{Threats to Validity}
\begin{itemize}
  \item \textbf{One machine class and synthetic data.} The suites are synthetic benchmark data with nearly fixed structure~\cite{duan2011}.
  \item \textbf{Unsupported cells.} QLever's two-order mode cannot answer patterns with a variable predicate; on suites that contain
  such patterns those cells are \emph{unsupported}, and the aligned ranking covers fewer suites.
  \item \textbf{Tuning.} We configured the engines; vendor tuning may differ. Jena is the answer reference, not a speed reference.
  \item \textbf{A null result.} If the rankings agree at $\tau \ge 0.8$ on every suite, the central claim fails and the paper reduces to the storage
  arm and the build results.
\end{itemize}

\section{Conclusion}
\todo{after Section 4.}

\appendix
\section{Artifact}
\todo{reproduction commands, archive layout \texttt{docs/eval-v1/}, SHA-256 of datasets.}

\begin{thebibliography}{9}
\bibitem{sidirourgos2008} L.~Sidirourgos, R.~Goncalves, M.~Kersten, N.~Nes, S.~Manegold. Column-store support for RDF data management: not all swans are white. \emph{PVLDB} 1(2), 2008.
\bibitem{duan2011} S.~Duan, A.~Kementsietsidis, K.~Srinivas, O.~Udrea. Apples and oranges: a comparison of RDF benchmarks and real RDF datasets. \emph{SIGMOD}, 2011.
\end{thebibliography}
\end{document}
